The professor requested a public GitHub-hosted site and API so that users can upload datasets and run the proposed methods. A practical deployment should separate the static site from the compute backend.

\section{Recommended Architecture}
\begin{itemize}
    \item GitHub repository: hosts source code, documentation, examples, and static frontend.
    \item Static frontend: GitHub Pages, allowing users to learn the method and upload datasets through a browser interface.
    \item API backend: FastAPI service deployed to a compute host such as Render, Railway, Hugging Face Spaces, or a university server. GitHub Pages alone cannot run Python APIs.
    \item Supported methods: M4, M6, and M7 as the three proposed methods; M1/M2 may be exposed as baselines.
    \item Supported input formats: JSON using the native fuzzy TOPSIS schema, flat CSV/XLSX with decision-maker, alternative, criterion, \(l,m,u\) columns, and crisp CSV/XLSX converted to pseudo-fuzzy ratings with user-specified criteria metadata.
\end{itemize}

\section{API Endpoints}
The proposed endpoints are:
\begin{itemize}
    \item \texttt{GET /health}: service status.
    \item \texttt{POST /api/validate}: validate uploaded dataset schema.
    \item \texttt{POST /api/run}: run M4, M6, and M7 and return rankings and closeness scores.
    \item \texttt{POST /api/convert}: convert flat CSV/XLSX into the native JSON schema.
    \item \texttt{GET /api/example}: return a sample dataset.
\end{itemize}

\section{Reproducibility}
The site should display method parameters, seed values, bag counts, and dataset metadata with every result. This is important for research transparency and reproducibility.
